\usepackage{ifthen}
\input{agenda/common.inc.tex}

% Title
\ifthenelse{\equal{\agendalanguage}{french}}{
  \def \trainingtitle{Formation sécurité des systèmes Linux embarqués}
}{
  \def \trainingtitle{Embedded Linux Security}
}

\def \trainingicon{common/flaticon-security-training.png}

% Duration
\ifthenelse{\equal{\sessiontype}{online}}{
  \def \trainingduration{4}
}{
  \def \trainingduration{3}
}

% Training objectives
\ifthenelse{\equal{\agendalanguage}{french}}{
  \def \traininggoals{
    \begin{itemize}
    \item Être capable de comprendre les mécanismes de sécurité et d'isolation des
      systèmes Linux embarqués modernes : mémoire non exécutable, randomisation
      de l'espace d'adressage, niveaux de privilèges, contrôle d'accès obligatoire
      et discrétionnaire.
    \item Être capable de concevoir et de mettre en œuvre une chaîne de démarrage
      sécurisée, du chargeur de démarrage jusqu'à l'espace utilisateur, y compris
      dm-verity.
    \item Être capable de raisonner sur la cryptographie et de mettre en œuvre un
      stockage sécurisé des clés, à l'aide d'un HSM matériel ou logiciel.
    \item Être capable d'exploiter la technologie ARM TrustZone (monde sécurisé,
      TF-A, OP-TEE) afin d'isoler les opérations sensibles et de développer des
      Trusted Applications.
    \item Être capable de mettre en œuvre des mécanismes de sécurité en espace
      utilisateur : capacités Linux, namespaces, cgroups, SECCOMP, SELinux et
      fonctionnalités de durcissement de systemd.
    \item Être capable de configurer le chiffrement du système de fichiers à
      l'aide de dm-crypt et de l'intégrer à une gestion sécurisée des clés via
      PKCS\#11.
    \item Être capable de gérer les vulnérabilités logicielles à l'aide de bases
      de données CVE, de générer et d'analyser des Software Bills of Materials
      (SBoM), et de comprendre les exigences de conformité, notamment le Cyber
      Resilience Act.
    \item Être capable de concevoir et de déployer des mécanismes de mise à jour
      sécurisés à l'aide de schémas de partitionnement A/B et d'outils tels que
      RAUC ou SWUpdate.
    \item Être capable de comprendre les concepts de measured boot et de mettre
      en œuvre la vérification de l'intégrité de la plateforme à l'aide de TPM
      et de IMA/EVM.
    \end{itemize}
  }
}{
  \def \traininggoals{
    \begin{itemize}
    \item Be able to understand the security and isolation mechanisms of
      modern embedded Linux systems: non-executable memory, address
      space randomization, privilege levels, mandatory and discretionary
      access control.
    \item Be able to design and implement a secure boot chain, from
      the bootloader all the way to userspace, including dm-verity.
    \item Be able to reason about cryptography, and implement secure key
      storage, using hardware or software HSM.
    \item Be able to leverage ARM TrustZone technology (secure world,
      TF-A, OP-TEE) to isolate sensitive operations and develop Trusted
      Applications.
    \item Be able to implement userland security mechanisms: Linux
      capabilities, namespaces, cgroups, SECCOMP, SELinux, and systemd
      hardening features.
    \item Be able to set up filesystem encryption using dm-crypt and
      integrate it with secure key management through PKCS\#11.
    \item Be able to manage software vulnerabilities using CVE databases,
      generate and analyze Software Bills of Materials (SBoM), and
      understand compliance requirements including the Cyber Resilience Act.
    \item Be able to design and deploy secure update mechanisms using
      A/B partitioning schemes and tools like RAUC or SWUpdate.
    \item Be able to understand measured boot concepts and implement
      platform integrity verification using TPMs and IMA/EVM.
    \end{itemize}
  }
}

% Training prerequisites
\def \trainingprerequisites{
  \begin{itemize}
    \prerequisitecommandline
    \prerequisiteembeddedlinux
    \prerequisiteenglish
  \end{itemize}
}

% Training audience
\ifthenelse{\equal{\agendalanguage}{french}}{
  \def \trainingaudience{
    Entreprises et ingénieurs qui conçoivent, développent et produisent
    des systèmes Linux embarqués ayant des exigences en matière de
    cybersécurité.
  }
}{
  \def \trainingaudience{
    Companies and engineers who design, develop and produce embedded
    Linux systems that have cyber-security requirements.
  }
}

\def \trainers {mathieu-dubois-briand,olivier-benjamin}

% Time ratio
\def \onsitelecturetimeratio{40}
\def \onsitelabtimeratio{60}

% Agenda items

\defagendaitem
{overview}
{lecture}
{Fundamental concepts}
{
  \begin{itemize}
  \item Threat modeling
  \item Cryptography basics: public/private key
  \item Understanding a TLS handshake from start to finish
  \item Understanding Public Key Infrastructures
  \end{itemize}
}
{Concepts fondamentaux}
{
  \begin{itemize}
  \item Modélisation des menaces
  \item Bases de la cryptographie : clé publique / clé privée
  \item Comprendre un handshake TLS du début à la fin
  \item Comprendre les infrastructures à clés publiques (PKI)
  \end{itemize}
}

\defagendaitem
{pki}
{lab}
{Setting up a small PKI}
{
  \begin{itemize}
  \item Using OpenSSL to generate a key hierarchy
  \item Experimenting with key revocation
  \item Performance comparison of public/private key encryption
  \end{itemize}
}
{Configuration d'une PKI simple}
{
  \begin{itemize}
  \item Utilisation d’OpenSSL pour générer une hiérarchie de clés
  \item Expérimentation de la révocation de clés
  \item Comparaison des performances du chiffrement par clé publique/clé privée
  \end{itemize}
}

\defagendaitem
{hwsec}
{lecture}
{Harware-enforced security barriers}
{
  \begin{itemize}
  \item General security concepts: kernel/user split, CPL, memory protections
  \item Common mitigations: ASLR, NX/DEP, SMAP/SMEP, RELRO
  \item ARM features: Exception levels
  \item ARM features: Secure world, TF-A, OP-TEE
  \item Understanding a full ARMv8 bootflow, including secure world
  \item Understanding Trusted Applications
  \end{itemize}
}
{Limites de sécurité assistées par le matériel}
{
  \begin{itemize}
\item Concepts généraux de sécurité : séparation noyau/utilisateur,
  CPL, protections mémoire
\item Contre-mesures courantes : ASLR, NX/DEP, SMAP/SMEP, RELRO
\item Fonctionnalités ARM : niveaux d’exception
\item Fonctionnalités ARM : {\em secure world}, TF-A, OP-TEE
\item Compréhension d’un flux de démarrage complet ARMv8, y compris le
  monde sécurisé
\item Compréhension des Trusted Applications
  \end{itemize}
}

\defagendaitem
{hwsec}
{lab}
{Exploring the secure world}
{
  \begin{itemize}
  \item Building secure world software (TF-A, OP-TEE, TA)
  \item Adding logs to observe Exception level transitions
  \item Adding logs to observe world transitions
  \item Writing a small Trusted Application in OP-TEE
  \item Interacting with our Trusted Application from userland
  \end{itemize}
}
{Exploration du {\em secure world}}
{
  \begin{itemize}
  \item Compilation de logiciel pour le {\em secure world} (TF-A,
    OP-TEE, TA)
  \item Ajout de journaux pour observer les transitions de niveaux
    d’exception
  \item Ajout de journaux pour observer les transitions entre mondes
  \item Écriture d’une petite Trusted Application dans OP-TEE
  \item Interaction avec notre Trusted Application depuis l’espace
    utilisateur
  \end{itemize}
}

\defagendaitem
{secureboot}
{lecture}
{Secure boot: part 1}
{
  \begin{itemize}
  \item Understanding the secure boot chain concept
  \item Secure boot implementation examples: x86 and ARM SoCs
  \item Secure boot: threat model
  \item Be able to design and implement a secure boot chain
  \end{itemize}
}
{Secure boot: première partie}
{
  \begin{itemize}
  \item Compréhension du concept de chaîne de démarrage sécurisée
  \item Exemples de mise en œuvre du secure boot : x86 et SoC ARM
  \item Secure boot : modèle de menace
  \item Être capable de concevoir et de mettre en œuvre une chaîne de
    démarrage sécurisée
  \end{itemize}
}

\defagendaitem
{secureboot}
{lab}
{Secure boot on i.MX93: ROM stage}
{
  \begin{itemize}
  \item Building all software components for secure boot (AHAB) on
    i.MX93
  \item Creating a signed AHAB container using SPSDK
  \item Flashing the fuses on the i.MX93 to enable secure boot (AHAB)
  \item Testing that secure boot (AHAB) is correctly enabled
  \end{itemize}
}
{Secure boot sur i.MX93}
{
  \begin{itemize}
  \item Compilation de tous les composants logiciels pour le secure
    boot (AHAB) sur i.MX93
  \item Création d’un conteneur AHAB signé à l’aide de SPSDK
  \item Programmation des fusibles sur l’i.MX93 pour activer le secure
    boot (AHAB)
  \item Vérification que le secure boot (AHAB) est correctement activé
  \end{itemize}
}

\defagendaitem
{secureuboot}
{lecture}
{Secure boot: part 2}
{
  \begin{itemize}
  \item Deep dive into the U-Boot components: TPL/SPL, U-Boot proper,
    FIT images
  \item Hardware description: Device Tree
  \item Enabling signature verification in the U-Boot bootloader
  \item RootFS verification: {\em dm-verity} and implications for system
    design
  \end{itemize}
}
{}
{
  \begin{itemize}
  \item Analyse approfondie des composants d’U-Boot: TPL/SPL, U-Boot principal, images FIT
  \item Description matérielle: Device Tree
  \item Activation de la vérification des signatures dans le chargeur de démarrage U-Boot
  \item Vérification du RootFS: {\em dm-verity} et implications pour la conception du système
  \end{itemize}
}

\defagendaitem
{secureuboot}
{lab}
{Secure boot on i.MX93: bootloader stage}
{
  \begin{itemize}
  \item Integrating the public key into the DTB
  \item Adding signature to the kernel FIT image
  \item Configuring U-Boot to enforce signature verification
  \item (optional) Configuring SPL to also enforce signature verification
  \end{itemize}
}
{Secure boot sur i.MX93: étape du bootloader}
{
  \begin{itemize}
  \item Intégration de la clé publique dans le DTB
  \item Ajout d’une signature à l’image FIT du noyau
  \item Configuration de U-Boot pour imposer la vérification des signatures
  \item (optionnel) Configuration du SPL pour imposer également la vérification des signatures
  \end{itemize}
}

\defagendaitem
{filesystem}
{lecture}
{Confidential information}
{
  \begin{itemize}
  \item Filesystem encryption, presentation of dm-crypt/cryptsetup
  \item Key management: Hardware Security Module
  \item Secure key usage: Introduction to PKCS\#11
  \item Secure key usage: kernel keyring and trusted keys
  \end{itemize}
}
{Protection des données}
{
  \begin{itemize}
  \item Chiffrement du système de fichiers, présentation de dm-crypt/cryptsetup
  \item Gestion des clés : module matériel de sécurité (HSM)
  \item Utilisation sécurisée des clés : introduction à PKCS\#11
  \item Utilisation sécurisée des clés : trousseau de clés du noyau et clés de confiance
  \end{itemize}
}

\defagendaitem
{keymanagement}
{lab}
{Secure key management}
{
  \begin{itemize}
  \item Provisioning keys into the i.MX93 ELE
  \item Using OP-TEE as software HSM
  \item (optional) Filesystem encryption using the ELE key
  \item (optional) Signing U-Boot FITs with HSM integration over PKCS\#11
  \end{itemize}
}
{Gestion des clés}
{
  \begin{itemize}
  \item Provisionnement des clés dans l’ELE du i.MX93
  \item Utilisation d’OP-TEE comme HSM logiciel
  \item (optionnel) Chiffrement du système de fichiers à l’aide de la clé ELE
  \item (optionnel) Signature des images FIT de U-Boot avec intégration d’un HSM via PKCS\#11
  \end{itemize}
}

\defagendaitem
{userland}
{lecture}
{Userland security measures}
{
  \begin{itemize}
  \item Access Control paradigms: MAC/DAC
  \item Linux capabilities: usage and examples
  \item Process isolation: Namespaces, Cgroups, SECCOMP
  \item Linux Security Modules: SELinux, AppArmor
  \item Application hardening via {\em systemd}
  \end{itemize}
}
{Mécanismes de sécurité en user-space}
{
  \begin{itemize}
  \item Paradigmes de contrôle d’accès : MAC/DAC
  \item Linux {\em capabilities} : utilisation et exemples
  \item Isolation des processus : namespaces, cgroups, SECCOMP
  \item Modules de sécurité Linux : SELinux, AppArmor
  \item {\em Hardening} des applications via {\em systemd}
  \end{itemize}
}

\defagendaitem
{userland}
{lab}
{Restricting userland applications}
{
  \begin{itemize}
  \item Preventing a simple application from executing shellcode using SECCOMP
  \item Manipulating SELinux contexts
  \item Using systemd to restrict a daemon's access to resources
  \end{itemize}
}
{Restriction des applications user-space}
{
  \begin{itemize}
  \item Empêcher une application simple d’exécuter du shellcode à l’aide de SECCOMP
  \item Manipulation des contextes SELinux
  \item Utilisation de systemd pour restreindre l’accès d’un démon aux ressources
  \end{itemize}
}

\defagendaitem
{measuredboot}
{lecture}
{Measured Boot}
{
  \begin{itemize}
  \item Boot event logs
  \item TPM's role in platform integrity
  \item Introduction to IMA/EVM
  \end{itemize}
}
{Measured Boot}
{
  \begin{itemize}
  \item Journaux des événements de démarrage
  \item Rôle du TPM dans l’intégrité de la plateforme
  \item Introduction à IMA/EVM
  \end{itemize}
}

\defagendaitem
{compliance}
{lecture}
{Maintenance, regulations and compliance}
{
  \begin{itemize}
  \item Introduction to vulnerability frameworks: CVE, CWE, CVSS
  \item A walkthrough of the Cyber Resilience Act (CRA)
  \item Defining an upgrade strategy, based on the release and
    maintenance cycles of important open-source projects
  \item Generating a Software Bill of Materials
  \item Analyzing a Software Bill of Materials
  \end{itemize}
}
{Maintenance, réglementations et conformité}
{
  \begin{itemize}
  \item Introduction aux cadres de gestion des vulnérabilités : CVE,
    CWE, CVSS
  \item Présentation du Cyber Resilience Act (CRA)
  \item Définition d’une stratégie de mise à jour, basée sur les
    cycles de publication et de maintenance des projets open source
    importants
  \item Génération d’un Software Bill of Materials
  \item Analyse d’un Software Bill of Materials
  \end{itemize}
}

\defagendaitem
{compliance}
{lab}
{Getting familiar with SBoMs and CVEs}
{
  \begin{itemize}
  \item Generating an SBoM for a Yocto distribution
  \item Analyzing the generated SBoM using \code{sbom-cve-check}
  \item Patching a CVE present in the SBoM
  \item Differential analysis
  \end{itemize}
}
{Prise en main des SBoM et des CVE}
{
  \begin{itemize}
  \item Génération d’un SBoM pour une distribution Yocto
  \item Analyse du SBoM généré à l’aide de \code{sbom-cve-check}
  \item Correction d’une CVE présente dans le SBoM
  \item Analyse différentielle
  \end{itemize}
}
\defagendaitem
{updates}
{lecture}
{Secure updates}
{
  \begin{itemize}
  \item A/B updates
  \item Presentation of SWUpdate
  \item Presentation of RAUC
  \item Bootloader integration
  \item Consequences on the secure boot chain
  \end{itemize}
}
{Mises à jour sécurisées}
{
  \begin{itemize}
  \item Mises à jour A/B
  \item Présentation de SWUpdate
  \item Présentation de RAUC
  \item Intégration au chargeur de démarrage
  \item Conséquences sur la chaîne de démarrage sécurisée
  \end{itemize}
}

\defagendaitem
{updates}
{lab}
{Setting up RAUC for secure updates}
{
  \begin{itemize}
  \item RAUC configuration on the target
  \item Building, shipping and installing a RAUC bundle 
  \item Bundle signature verification using PKCS\#11
  \item (optional) RAUC bundle encryption
  \item (bonus) HSM integration
  \end{itemize}
}
{Mise en oeuvre de RAUC pour les mises à jour sécurisées}
{
  \begin{itemize}
  \item Configuration de RAUC sur la cible
  \item Création, déploiement et installation d’un bundle RAUC
  \item Vérification de la signature du bundle à l’aide de PKCS\#11
  \item (optionnel) Chiffrement du bundle RAUC
  \item (bonus) Intégration d’un HSM
  \end{itemize}
}

\def \agendaboards{imx93-frdm}

\def \onsiteagenda {
  \ifthenelse{\equal{\agendalanguage}{french}}{
    \section{Programme de la formation}
  }{
    \section{Training Schedule}
  }
  \begin{tabularx}{\textwidth}{p{2cm}p{5cm}p{11cm}}
  \showagendaday{1}
  \showagendaitem{overview}{lecture}
  \showagendaitem{pki}{lab}
  \showagendaitem{hwsec}{lecture}
  \showagendaday{2}
  \showagendaitem{hwsec}{lab}
  \showagendaitem{secureboot}{lecture}
  \showagendaitem{secureboot}{lab}
  \showagendaday{3}
  \showagendaitem{secureuboot}{lecture}
  \showagendaitem{secureuboot}{lab}
  \showagendaday{4}
  \showagendaitem{filesystem}{lecture}
  \showagendaitem{keymanagement}{lab}
  \showagendaitem{userland}{lecture}
  \showagendaday{5}
  \showagendaitem{userland}{lab}
  \showagendaitem{measuredboot}{lecture}
  \showagendaitem{compliance}{lecture}
  \showagendaitem{compliance}{lab}
  \showagendaday{6}
  \showagendaitem{updates}{lecture}
  \showagendaitem{updates}{lab}
  \end{tabularx}
}

\def \onlineagenda {
  \ifthenelse{\equal{\agendalanguage}{french}}{
    \section{Programme de la formation}
  }{
    \section{Training Schedule}
  }
  \begin{tabularx}{\textwidth}{p{2cm}p{5cm}p{11cm}}
  \showagendaday{1}
  \showagendaitem{overview}{lecture}
  \showagendaitem{pki}{lab}
  \showagendaitem{hwsec}{lecture}
  \showagendaitem{hwsec}{lab}
  \showagendaday{2}
  \showagendaitem{secureboot}{lecture}
  \showagendaitem{secureboot}{lab}
  \showagendaitem{secureuboot}{lecture}
  \showagendaitem{secureuboot}{lab}
  \showagendaday{3}
  \showagendaitem{filesystem}{lecture}
  \showagendaitem{keymanagement}{lab}
  \showagendaitem{userland}{lecture}
  \showagendaitem{userland}{lab}
  \showagendaday{4}
  \showagendaitem{compliance}{lecture}
  \showagendaitem{compliance}{lab}
  \showagendaitem{updates}{lecture}
  \showagendaitem{updates}{lab}
  \end{tabularx}
}
